% Input work/deg10/next_cases_20261030/dyadic_formal_kernel.tex before this file.
\subsection{Excluding a point of order $49$}
\begin{proposition}\label{prop:decic-no49}
An elliptic curve over a number field of degree $10$ has no point of order $49$.
\end{proposition}
\begin{proof}
Put $X=X_1(49)$ and $J=J_1(49)$. By
\cite[Theorems~3.1 and 4.13, and Table~2]{DEHMZ},
\[
 J(\mathbb Q)\simeq C_7\times C_{52367710906884085342}.
\]
In particular $J(\mathbb Q)[2^\infty]\simeq C_2$.
The map $\pi:X\to X_0(49)$ has degree $21$ and diamond group
$H\simeq C_{21}$. The latter acts trivially on the rational
$2$-primary subgroup, so
$\pi^*\pi_*t=\sum_{h\in H}ht=21t$ makes $\pi_*$ injective there.
Similarly $\pi_*\pi^*=[21]$ makes pullback injective on
$J_0(49)(\mathbb Q)[2^\infty]$, whose order is therefore at most two.
The model
\[
 y^2+xy=x^3-x^2-2x-1
\]
for $J_0(49)$ has discriminant $-343$ and a point of order two
$(2,-1)$, reducing to the nonzero point $(0,1)$ modulo $2$.
Thus reduction is injective on the rational $2$-primary subgroup
of $J_0(49)$, hence also on that of $J$. It is always injective on
odd-order torsion. Therefore
\[
 J(\mathbb Q)\longrightarrow J(\mathbb F_2)
\]
is injective.

Suppose that $D$ is a rational effective noncuspidal divisor of
degree $10$ on $X$. Exhaustive enumeration over $\mathbb F_{2^f}$,
$f\leq10$, shows that the only possible noncuspidal closed-point
degrees are $9$ and $10$. There are respectively $42$ and $126$
such points, forming two and six diamond orbits. For every orbit
representative $P$, the degree-zero divisor
\[
 A_3P=(T_3-1-3\langle3\rangle)P
\]
is nonprincipal, by exact Riemann--Roch computations on the full
genus-$69$ curve. Independent enumeration of all curves, twists
and level structures verifies the completeness of these orbits.
The four native $3$-isogenies at each of the eight representatives
also reproduce the finite correspondence, including all closed
point degrees and norm multiplicities.

The operator $A_3$ kills $J(\mathbb Q)$, by Eichler--Shimura at $3$
and injectivity of torsion reduction there. At a rational cusp $c$
we have the literal identity $T_3c=c+3\langle3\rangle c$.
It follows from the Tate parametrization at one rational cusp and
then from diamond equivariance at all $21$ rational cusps.
If the reduction of $D$ has noncuspidal support, it is therefore
one degree-$10$ point or a degree-$9$ point plus a rational cusp.
In either case the computed nonprincipality of $A_3P$ is a
contradiction.

The cusp degrees over $\mathbb Q$ and $\mathbb F_2$ are
\[
 1^{21},6^3,21\qquad\text{and}\qquad1^{21},3^6,21.
\]
If the reduction of $D$ is supported on rational cusps, let $C$
be its rational cuspidal lift. Injectivity of specialization gives
$D\sim C$. Since
\[
 \operatorname{gon}_{\mathbb C}X
 \geq \frac7{800}[\operatorname{PSL}_2(\mathbb Z):\overline{\Gamma_1(49)}]
 =\frac7{800}\,1176=\frac{1029}{100}>10,
\]
we obtain $D=C$, a contradiction.

It remains to exclude cuspidal reductions containing a cubic cusp.
Put $U=(\mathcal C-1)H^0(X_{\mathbb F_2},\Omega^1)$, where
$\mathcal C$ is Cartier. Its dimension is $68$. We verify that
restriction from $U$ has rank $10$ at every effective cuspidal
divisor of degree $10$ containing a cubic cusp, as follows.

First use the existing common-basis jets at eighteen rational
cusps, through multiplicity six, and at all six cubic cusps,
through multiplicity three. The complete degree-$10$ replay checks
\[
 6\left(\binom{24}{7}-18\right)
 +\binom72\binom{21}{4}+\binom83\,18
 =2203209
\]
divisors and finds rank $10$ in every case. This covers every
mixed cuspidal divisor having at most six rational support points
and rational multiplicities at most six: the diamond group acts
regularly on the $21$ rational cusps, and at most $6\cdot3=18$
of its $21$ elements send some support point to one of the three
omitted cusps. Cartier commutes with diamonds.

If there are seven rational support points, the divisor consists
of seven distinct rational cusps and one cubic cusp. A rational
multiplicity of seven instead forces the divisor to be $7c+Q$,
where $Q$ is cubic. Diamond transitivity fixes the choice of $c$,
so the latter case requires six checks.

We obtain the remaining jets as follows. Let $\mathsf C$ be the
$69\times69$ Cartier matrix in the existing differential basis.
An exact native calculation of the pullback by $\langle2\rangle$
gives a matrix $\mathsf D$. It has order $21$, commutes with
$\mathsf C$, and belongs to $\mathbb F_2[\mathsf C]$.
The minimal polynomial of $\mathsf C$ has degree $69$.
Thus the native matrix $\mathsf D$ remains valid under any change
of basis preserving the Cartier matrix. Indeed, such a change
commutes with every polynomial in $\mathsf C$.
The orbit under $\mathsf D$ of one rational-cusp evaluation
column $a_0$, based at a cusp $c_0$, consists of $21$ distinct columns and contains all
$18$ previously recorded rational-cusp columns. It therefore
recovers the three missing first-order restrictions. The verifier
also checks the covariance of all six cubic restriction subspaces
through orders one, two and three.

For the seventh-order restriction at a rational cusp we use a
fresh native jet. Its Cartier matrix agrees with $\mathsf C$,
and its first-order column is $a_0$. At a rational cusp over
$\mathbb F_2$, the nonzero cotangent trivialization is unique. The columns
$a_0,\mathsf C a_0,\ldots,\mathsf C^{68}a_0$ span all $69$
dimensions. This verifies the required basis compatibility,
without an assumption on the ordering of the native cusps.
To see this, let $S$ be the change of differential basis and write
the fresh base cusp as $\langle2\rangle^h c_0$. Then
$S\mathsf C=\mathsf C S$ and
\[
 a_0=S\mathsf D^h a_0.
\]
Since $\mathsf D$ is a polynomial in $\mathsf C$ and $a_0$ is
cyclic for $\mathsf C$, we obtain $S\mathsf D^h=I$.
Thus $S$ is a diamond power. It permutes the six cubic restriction
subspaces, while changes of local parameters and differential
units apply invertible operations to the restriction matrices.
Testing all six cubic cusps therefore makes the fresh
seventh-order jet compatible with the existing cubic data.
Magma's built-in local completion independently verifies all seven
rational-cusp restriction kernels.

The independent verifier checks all
$6\binom{21}{7}=697680$ seven-distinct-rational-cusp divisors and
all six divisors $7c_0+Q$. Every restriction has rank $10$.
It also verifies $\mathsf D\in\mathbb F_2[\mathsf C]$, the exact
order and Cartier commutation, all $18$ known rational evaluators,
and all cubic covariance identities. Hence every mixed cuspidal
divisor of degree $10$ has the required restriction rank.

Corollary~\ref{cor:dyadic-Cartier-formal-immersion}, applied with the full Jacobian as
target, now excludes the remaining cuspidal reductions. Each
closed cubic cusp is lifted with its complete Frobenius orbit,
so the entire cuspidal lift is a $\mathbb Q_2$-rational divisor.
This completes the proof.
\end{proof}
